\chapter{Oscillatore Armonico: soluzione analitica}

Consideriamo un oscillatore armonico quantistico unidimensionale, il cui hamiltoniano è esprimibile nella forma 
% 
\begin{equation}
    \hat{H} = \frac{\hat{P}^2}{2m} + \frac{m\omega^2 \hat{Q}^2}{2}
\end{equation}
Possiamo trovare lo spettro degli autovalori proseguendo con due approcci diversi: con un \emph{approccio algebrico} o con un \emph{approccio analitico}. 

In \cite{Ballentine} sono affrontati entrambi gli approcci, tuttavia esiste un modo che semplifica notevolmente l'approccio analitico che permette di avere a che fare solo con equazioni differenziali del primo ordine (mentre altrimenti dovremmo trattare equazioni differenziali di secondo ordine). \\
Questo metodo parte dai risultati ottenuti con l'approccio algebrico, per questo sarà necessario riportarli velocemente. 

Definendo l'operatore di distruzione\mkeyword{operatore di distruzione} e l'operatore di costruzione\mkeyword{operatore di costruzione} rispettivamente come 
%
\begin{equation}
\label{eq:cost_distr}
    \hat{a} = \frac{1}{\sqrt{2}} \left( \frac{\hat{Q}}{l_{\omega}} + i \frac{\hat{P}}{p_{\omega}} \right) \quad , \quad \hat{a}^{\dagger} = \frac{1}{\sqrt{2}} \left( \frac{\hat{Q}}{l_{\omega}} - i \frac{\hat{P}}{p_{\omega}} \right)
\end{equation}
E definendo $\{\ket{n} : \ n\in \mathbb{N} \cup \{ 0 \}\}$ l'insieme degli autostati di $\hat{H}$ si ha che valgono le seguenti relazioni
%
\begin{align}
    \hat{a} \ket{0} &= 0 \label{eq:prima_rel}\\
    \ket{n} &= \frac{1}{\sqrt{n!}}(\hat{a}^{\dagger})^n \ket{0}
    \label{eq:relazione_creaz_distr}
\end{align}
La prima servirà per trovare $\psi_0 (x)$ mentre la seconda sarà necessaria per ricavare $\psi_n(x)$. 

Vale infatti
%
\begin{align}
    \Braket{x|\hat{a}|0} &= \frac{1}{\sqrt{2}}\Braket{x|\frac{\hat{Q}}{l_{\omega}} + i l_{\omega} \frac{\hat{P}}{\hbar}|0} \\
    &= \frac{1}{\sqrt{2}}\left( \frac{x}{l_{\omega}} + l_{\omega} \partial_x \right) \psi_0(x) = 0
\end{align}
Dove abbiamo usato la prima relazione, riscrivendo $\hat{Q}$ e $\hat{P}$ in coordinate. Inoltre abbiamo usato $l_{\omega} \cdot p_{\omega} = \hbar$.

Risolvendo l'equazione differenziale con il metodo delle variabili separabili si ottiene 
%
\begin{equation*}
    \psi_0(x) = A e^{- \frac{x^2}{2 l_{\omega}^2}}
\end{equation*}
Il valore di $A$ si trova con la normalizzazione 
%
\begin{align*}
    \int_{-\infty}^{\infty} dx \ |\psi_0(x)|^2 = A^2 \int_{-\infty}^{\infty} dx \ e^{-\frac{x^2}{l_{\omega}^2}} = A^2 \ \sqrt{\pi l_{\omega}^2} = 1 \implies A = \frac{1}{l_{\omega}^{1/2}\pi^{1/4}}
\end{align*}
%
Quindi si ha 
%
\begin{equation}
\label{eq:psi_0}
    \psi_0(x) = \frac{1}{l_{\omega}^{1/2}\pi^{1/4}} e^{- \frac{x^2}{2 l_{\omega}^2}}
\end{equation}

Adesso grazie alla relazione \eqref{eq:relazione_creaz_distr} possiamo ricavare $\psi_n(x) \ , \forall n \ \in \mathbb{N} \cup \{ 0 \}$. Ricaveremo solo $\psi_1(x)$ per poi dare una formula generale per $\psi_n(x)$, tuttavia sarebbe possibile ricavare ogni $\psi_n(x)$ con un metodo del tutto analogo. 
%
\begin{align}
    \psi_1(x) = \Braket{x|1} = \Braket{x|\hat{a}^{\dagger}|0} &= \frac{1}{\sqrt{2}}\Braket{x|\frac{\hat{Q}}{l_{\omega}} - i l_{\omega} \frac{\hat{P}}{\hbar}|0}\\
    &= \frac{1}{\sqrt{2}} \left( \frac{x}{l_{\omega}} - l_{\omega} \partial_x \right) \psi_0(x)
\end{align}
Adesso sostituendo \eqref{eq:psi_0} e derivando si ottiene 
%
\begin{equation}
    \psi_1(x) = \left( \frac{2}{\pi^{1/2}l_{\omega}^3} \right)^{1/2} x \ e^{-\frac{x^2}{2 l_{\omega}^2}}
\end{equation}
Potremmo svolgere un processo simile per ricavare $\psi_n(x)$, una formula generale è però 
%
\begin{tcolorbox}[colback=black!5!white, colframe=black!45!white, boxrule=0.8pt]
\begin{equation}
    \psi_n(x) = \left( \frac{1}{\pi^{1/2} l_{\omega} 2^n n!} \right)^{1/2} H_n\left(\frac{x}{l_{\omega}}\right) \ \text{exp}\left( -\frac{x^2}{2 l_{\omega}^2} \right)
\end{equation}
\end{tcolorbox}
Dove si sono introdotti i polinomi di Hermite\mkeyword{polinomi di Hermite}:
\begin{equation}
    H_n(y) = e^{y^2/2} \left( y - \frac{d}{dy} \right)^n e^{-y^2/2}
\end{equation}

\section{Base degli impulsi}
Nei passaggi precedenti abbiamo deciso di trovare gli autostati dell'hamiltoniano in coordinate, tuttavia possiamo farlo in modo del tutto analogo anche usando come base gli autostati dell'operatore impulso. 

Sfruttiamo, come prima, la \eqref{eq:prima_rel} per ottenere $\tilde{\psi}_0(p)$, l'unica cosa è che adesso usiamo $\bra{p}$ invece di $\bra{x}$
%
\begin{align}
    \Braket{p|\hat{a}|0} &= \frac{1}{\sqrt{2}}\Braket{p|p_{\omega}\frac{\hat{Q}}{\hbar} + i \frac{\hat{P}}{p_{\omega}}|0} \\
    &= \frac{1}{\sqrt{2}}\left( i p_{\omega} \partial_p +  i \frac{p}{p_{\omega}}\right) \tilde{\psi}_0(p) = 0
\end{align}
Risolvendo l'equazione differenziale e imponendo la normalizzazione si ha 
\begin{equation}
    \tilde{\psi}_0(p) = \frac{1}{p_{\omega}^{1/2} \pi^{1/4}} e^{-\frac{p^2}{2 p_{\omega}^2}}
\end{equation}
Allora è facile intuire che anche le altre funzioni d'onda saranno analoghe al caso precedente e basterà effettuare lo scambio $x \to p$ e $l_{\omega} \to p_{\omega}$
\begin{equation}
    \tilde{\psi}_n(p) = \left( \frac{1}{\pi^{1/2} p_{\omega} 2^n n!} \right)^{1/2} H_n\left(\frac{p}{p_{\omega}}\right) \ \text{exp}\left( -\frac{p^2}{2 p_{\omega}^2} \right)
\end{equation}